\section{Boolean完全交叉上の全一基点二次式のフレーム}
\label{sec:boolean-quadric-boundary}
前節では最後の二次式を$e_3$で割り切れるものに限った。
ここでは交差項も許し、最初の三式を固定した場合の一基点二次式を全て扱うフレームを与える。
これ自体からAS正則性や非Koszul性は結論しない。

\begin{proposition}[全一基点開集合での三次評価フレーム]
\label{prop:boolean-global-cubic-frame}
$Q=(e_0^2-e_0e_3,e_1^2-e_1e_3,e_2^2-e_2e_3)$とし、
$\epsilon\in\{0,1\}^3$に対応する列ベクトルを$p_\epsilon=(\epsilon_0,\epsilon_1,\epsilon_2,1)$とする。
二次式$q$が$p_{000}$でのみ零となり、他の七点で非零であるとする。
$Q$の線型結合を加え、非零スカラー倍で正規化すれば、$q$は一意に
\[
 q=e_0e_3+u e_1e_3+v e_2e_3
 +(a-1-u)e_0e_1+(b-1-v)e_0e_2+(c-u-v)e_1e_2
\]
と書ける。係数域は
\[
 X=\{uvabc(a+b+c-1-u-v)\ne0\}\subset\mathbb A^5
\]
であり、七点での値は順に$1,u,v,a,b,c,a+b+c-1-u-v$である。
$B=\operatorname{Sym}(V^*)/(Q,q)$に対し、
\[
 b_i=\sum_{\epsilon_i=1}
       \frac{(-1)^{|\epsilon|-1}}{q(p_\epsilon)}p_\epsilon^{\otimes3}
       \quad(i=0,1,2),\qquad b_3=p_{000}^{\otimes3}
\]
は$X$の全域で正則な$B_3^*$の基底を与える。
さらに
\[
 (\dim B_n)_{n\geq0}=(1,4,6,4,1,1,1,\ldots)
\]
が全ての$X$上で成立する。
\end{proposition}
\begin{proof}
$Q$を使って$e_i^2$を$e_ie_3$へ置換すると、二次式の標準形は
$e_3^2,e_ie_3,e_ie_j$ $(i<j<3)$の七項である。
$p_{000}$での零条件が$e_3^2$係数を消し、$q(p_{100})\ne0$が$e_0e_3$係数の正規化を許す。
残りは上記の五係数と七点での評価から一意に決まる。

$R=\operatorname{Sym}(V^*)/Q$の次数3以上の成分は、八点の評価によって$k^8$と同型になる。
実際、次数3の多項式
$\prod_{i:\epsilon_i=1}e_i\prod_{i:\epsilon_i=0}(e_3-e_i)$
が各点のindicatorを与える。
従って八つの三次評価テンソル$p_\epsilon^{\otimes3}$は独立である。
$R_1$の七非零点への評価は階数4なので、$qR_1\subset R_3$も4次元である。
よって$\dim B_3=8-4=4$となる。

各$b_i$は$Q$を消す。$q$と任意の一次式との積の評価では分母が消え、
$\epsilon_i=1$の面の四頂点の符号付き和が零になるので、$qR_1$も消す。
$b_3$も$q(p_{000})=0$からこれを消す。
四つの評価重み列を$p_{100},p_{010},p_{001},p_{000}$の行で見ると、
行列は$\operatorname{diag}(1,u^{-1},v^{-1},1)$となる。
従って四テンソルは全ての$X$上で独立であり、基底をなす。

最後に$R_2$は7次元で、七非零点への評価は同型である。
もし次数2の元がその七点で消えるなら、その八点上の関数は$p_{000}$のindicatorの定数倍である。
しかしその一意なmultilinear表示は次数3の項を持つので、次数2では非零な倍を表せない。
従って$\dim B_4=8-7=1$となる。
$n\geq5$では$R_{n-2}$の全八点評価が同型であり、$q$の乗法の階数は7なので、$\dim B_n=1$である。
低次数の値は四つの独立な二次式から従う。
\end{proof}

\begin{corollary}[五係数域全体での対称成分の閉包への所属]
\label{cor:boolean-family-in-symmetric-closure}
前命題の基底から$(w_q)_{abjc}=(b_j)_{abc}$と定めたテンソルは、全て$Z_{\mathrm{sym}}$に属する。
\end{corollary}
\begin{proof}
$q_t=q+t e_3^2$とする。
$i=0,1,2$の基底式では分母を$q(p_\epsilon)+t$へ置き換える。
第3基底は
\[
 b_3(t)=p_{000}^{\otimes3}
 -\frac{t}{u+t}p_{010}^{\otimes3}
 -\frac{t}{v+t}p_{001}^{\otimes3}
 +\frac{t}{c+t}p_{011}^{\otimes3}
\]
とする。これは$\epsilon_0=0$の面の符号付き評価関係を$t$倍したものであり、
全て$q_tR_1$を消す。全ての分母は$t=0$で可逆であり、上の四基底へ特殊化する。
一般の$t$では$q_t$は全八点で非零なので、四二次式は基点自由完全交叉となる。
従って一般テンソルは対称チャートに属し、特殊値$w_q$はその飽和閉包に属する。
\end{proof}

\begin{remark}
このフレームは、$012,133,233,333$を指定する局所正規化の追加分母が零になっても定まる。
その局所正規化を一般点の計算に使う場合も、最後の特殊化議論は上の大域フレーム上で行う必要がある。
最初の三二次式を任意に動かした一般の一基点四二次式系まで被覆したとは主張しない。
\end{remark}

\begin{corollary}[Frobenius境界上の非対称化性]
\label{cor:boolean-frobenius-not-symmetric}
上の$X$で$w_q$が有限Frobenius条件を通るなら、$w_q\notin S$である。
従って、この係数域からKoszul AS点が得られれば、対称チャート外の低階数対象となる。
その存在をここで結論しているわけではない。
\end{corollary}
\begin{proof}
因子$0,1,3$は既に対称なので、$S_1,S_3$は$I$を含む。
対称化可能なら線型判定により$S_2$にも可逆元があり、それだけを因子2へ作用させれば完全対称になる。
この作用は因子$0,1$の縮約像$R_0$を変えない。
$R_0$は6次元であり、その直交二次関係は$\bigwedge^2(V^*)+(Q,q)$である。
従って得られた完全対称Frobenius対象の可換二次双対は$B$そのものとなる。
対称Frobeniusチャート定理はこれが基点自由完全交叉であることを要求するが、
$(Q,q)$には$p_{000}$が共通基点として残る。これは矛盾である。
\end{proof}

\begin{theorem}[固定Boolean完全交叉上の全一基点係数域のAS排除]
\label{thm:boolean-one-point-uniform-exclusion}
命題\ref{prop:boolean-global-cubic-frame}の全係数域$X$は、標準型Koszul AS点を一つも含まない。
有限Frobenius条件と全四次検査を通る任意の$X$上の点で、始点0の五次$H_2$は7次元以上である。
一般点では五次の$H_2$次元は$(7,0,0,0)$であり、$H_3$は全始点で零である。
この結論は最初の三二次式を固定した場合の全域を覆うが、任意の相異なる八点完全交叉への一様排除は主張しない。
\end{theorem}
\begin{proof}
各面の基底$b_i$を、その面上の四つの$q$評価値の積でスカラー倍すると、
$X$上の単元による基底変更で分母のない多項式テンソルを得る。
これを一般点の計算に使う。三因子対称性と四二次式の消去から、全二対二階数は6以下である。
$(u,v,a,b,c)=(1,1,2,2,2)$は前節の検算済み標本の三次基底変更であり、
その変更行列の行列式は$1/12$である。
従って有限Frobenius条件と全四次検査を通る部分$U\subset X$は空でない開部分である。

体$\mathbb Q(u,v,a,b,c)$上で商写像と自由単項式sectionを作り、
始点0の五次微分$d_2$を構成すると、正確な階数は77となる。
行列の大きさは$140\times120$である。
再現コード\path{computations/verify_boolean_quadric_homology.py}は、
浮動小数や有限個の評価点を用いず、五変数の有理関数の消去でこの値を検証する。
計算量を抑えるためpivotの順序を多項式の項数等で選ぶが、全ての操作は体上の可逆な消去である。

$U$上では$\operatorname{rk}d_3\leq36$であり、検算済みの点で36を達成する。
従って一般点での$d_3$階数は36であり、一般$H_2$次元は$120-77-36=7$である。
始点1,2,3の五次正確性と全始点の$H_3=0$は、同じ標本から開性によって一般点へ移る。
補題\ref{lem:generic-quintic-exclusion}を$X$と$U$へ適用すると、
全ての$U$上で$\dim H_2\geq7$となる。
$X\setminus U$では既に必要なFrobenius条件か四次検査を失っているので、そこにもAS点はない。
大域評価フレームを使っているため、途中の局所正規化の追加分母が消える値もこの議論に含まれる。
\end{proof}
